% Unaltered test distributions (docs/PREREGISTRATION_NATURAL.md; operating points: docs/PREREGISTRATION_REVIEW2.md, R3).
\section{Consequences on unaltered test data}\label{sec:natural}\label{sec:robust}
The traps construct an extreme reversal. Without any resampling, in-ROI artifacts are still associated with the label
in every cohort (thyroid calipers: 37.6\% of benign vs 24.1\% of malignant nodules; in-lesion hair: 31\% of melanomas
vs 20\% of benign lesions; capsule debris on the lesion: 23\% of polyp-like lesions vs 10\% of erosions). We trained on
natural training data and tested on untouched data: the official TN3K/TNCD test split, which is patient-disjoint from
the training split \citep{gong2022acl}, each ISIC source held out in turn (a real change of hospital and device), and
held-out capsule frame blocks.

\paragraph{Discrimination} All analyses in this section use the rebuilt runs (Sec.~\ref{sec:repro}).
On the shortcut-conflicting pairs masking \emph{lowers} AUROC in 3 of 5 test sets (thyroid $-0.095$ [$-0.167, -0.023$]; held-out BCN $-0.025$ [$-0.042, -0.008$]; held-out MSK $-0.065$ [$-0.111, -0.020$]), raises it for held-out HAM ($+0.050$ [$+0.030, +0.070$]), and does not change it significantly for capsule ($+0.021$ [$-0.003, +0.048$]). Over all test pairs masking changes thyroid AUROC by only $-0.004$ [$-0.048, +0.041$]: the harm is confined to the cases in which the in-ROI artifact points to the wrong diagnosis, which is exactly where a shortcut is dangerous%
.

\paragraph{Operating points (pre-registered)} Because a clinical decision uses a threshold, we fixed thresholds on
each seed's validation split at five operating points and compared test sensitivity and specificity
(Table~\ref{tab:op}; the bootstrap re-estimates the thresholds in every replicate). On the official thyroid split
masking lowers sensitivity at all five operating points, each with a CI excluding zero, for example from
0.699 to 0.424 at the maximal-balanced-accuracy threshold ($-0.275$ [$-0.399, -0.182$]) and from 0.557 to 0.355 at a validation specificity of 0.80 ($-0.202$ [$-0.265, -0.134$])%
. Specificity rises correspondingly, so part of the change is a shift of the score distribution relative to
thresholds learned on validation data (quantified below); the loss is largest where the shortcut conflicts with the
label: among the malignant nodules that carry an in-ROI caliper, at a validation specificity of 0.80 sensitivity falls
by $-0.372$ [$-0.436, -0.291$] (0.549 $\to$ 0.177)%
. \emph{Registration status for this frozen model: a pre-registered test inside a
post hoc subgroup.} The subgroup (the 78 malignant nodules of the official test split that carry an in-ROI caliper; the
same 78 in every seed) was defined after the first unaltered-data analysis, where its AUROC harm was found; the
operating-point test in it was registered before the operating points were computed, and for the fine-tuned model below
the whole test was registered before any result. The pre-registered decision rule
(loss at the maximal-balanced-accuracy threshold and at $\ge3$ of the four other operating points) is met. On capsule
endoscopy, where masking improves AUROC overall, it lowers the sensitivity for erosions carrying debris on the lesion
at the two high-sensitivity operating points (Supplement~S9).

\paragraph{A fine-tuned model near the benchmark (pre-registered)} To test whether these are properties of a weak
model, we compared twelve fine-tuning recipes on validation images only, committed the choice (ConvNeXt-T pretrained on
ImageNet-22k, learning rate $10^{-4}$, 8 epochs, 224\,px) and then trained it with five seeds. Its test AUROC is 0.769,
0.004 below the lowest published cross-entropy baseline on this split (0.773; \citealp{gong2022acl}) and above the
frozen probe (0.735). All four registered tests are supported (Holm): masking lowers the AUROC on shortcut-conflicting
pairs by $-0.102$ [$-0.178, -0.026$]; it lowers sensitivity at all four of OP2--OP5; at OP2 it lowers sensitivity among
the 78 malignant nodules with an in-ROI caliper from 0.744 to 0.223 ($-0.521$ [$-0.642, -0.382$]), a subgroup registered
before any result of this model existed; and its trap crossover is $+0.187$ [$+0.111, +0.268$]. Group balancing recovers
the conflicting pairs ($+0.088$ [$+0.019, +0.158$] over masking).

\paragraph{Is the loss a threshold artefact? (pre-registered)} Re-tuning each model's threshold on the test set
itself---the best that recalibration at a deployment site can achieve---makes the overall sensitivity of the two models
equal at a test specificity of 0.80 ($-0.017$ [$-0.124, +0.070$] for the fine-tuned model), but the conflicting subgroup
still loses: from 0.544 to 0.390 ($-0.154$ [$-0.302, -0.023$]) for the fine-tuned model and from 0.505 to 0.331
($-0.174$ [$-0.313, -0.012$]) for the frozen probe (Holm $p=0.019$ each), while the malignant nodules whose caliper
points the right way tend to gain. About 30\% of the loss at validation-fixed thresholds therefore cannot be removed by
any threshold; the rest follows from transporting thresholds from validation to test data, as deployment does. At a
matched specificity of 0.90 and a matched sensitivity of 0.80 the direction is the same, with intervals that include
zero (Supplement~S16).

\paragraph{New patients at natural prevalence (pre-registered)} Trained on all ISIC 2019 images and tested on ISIC 2020
without any resampling, masking lowers the AUROC on shortcut-conflicting pairs: $-0.026$ [$-0.052, -0.000$] under the
registered near-duplicate rule, which proved too loose (it removed most look-alike images of other patients), and
$-0.037$ [$-0.061, -0.012$] under a label-free calibrated rule registered as a sensitivity analysis after that result
(32,952 images). The overall AUROC is nearly unchanged ($+0.016$ [$-0.000, +0.033$]), and group balancing recovers the
conflicting pairs ($+0.072$ [$+0.048, +0.097$] over masking). No sensitivity loss at
a fixed operating point was detectable, and the cleanest in-lesion comparison (53 hair-free melanomas) was
inconclusive ($-0.005$ [$-0.057, +0.048$]): the external harm combines the removal of hair beside melanomas with the
retention of hair on benign lesions (Supplement~S16).
\begin{table}[t]\centering\scriptsize
\caption{Masking versus ERM on the unaltered, patient-disjoint official thyroid test split (614 images, 236 malignant) at five operating points fixed on validation data (mean over five seeds; $\Delta$ = mask $-$ ERM with 95\% crossed seed$\times$image bootstrap CI, thresholds re-estimated in every replicate). \emph{Conflicting}: the 78 malignant nodules with an in-ROI caliper (the caliper marks benign nodules in this dataset). The subgroup was defined after the first unaltered-data analysis; the operating-point test in it was registered before these numbers were computed.}\label{tab:op}
\resizebox{\linewidth}{!}{\begin{tabular}{lccccc}\toprule
Threshold on validation & Sens.\ ERM $\to$ mask & $\Delta$ sensitivity & Sens.\ conflicting ($n=78$) & $\Delta$ sensitivity, conflicting & $\Delta$ specificity \\\midrule
max.\ balanced accuracy & 0.699 $\to$ 0.424 & $-0.275$ [$-0.435, -0.138$] & 0.703 $\to$ 0.223 & $-0.479$ [$-0.640, -0.313$] & $+0.202$ [$+0.115, +0.330$] \\
specificity $\ge0.80$ & 0.557 $\to$ 0.355 & $-0.202$ [$-0.300, -0.100$] & 0.549 $\to$ 0.177 & $-0.372$ [$-0.499, -0.236$] & $+0.117$ [$+0.059, +0.172$] \\
specificity $\ge0.90$ & 0.371 $\to$ 0.219 & $-0.152$ [$-0.253, -0.043$] & 0.336 $\to$ 0.100 & $-0.236$ [$-0.383, -0.089$] & $+0.075$ [$+0.029, +0.122$] \\
sensitivity $\ge0.80$ & 0.753 $\to$ 0.380 & $-0.373$ [$-0.469, -0.270$] & 0.744 $\to$ 0.197 & $-0.546$ [$-0.671, -0.421$] & $+0.278$ [$+0.203, +0.359$] \\
sensitivity $\ge0.90$ & 0.869 $\to$ 0.588 & $-0.281$ [$-0.391, -0.178$] & 0.862 $\to$ 0.400 & $-0.462$ [$-0.614, -0.310$] & $+0.339$ [$+0.236, +0.425$] \\
\bottomrule\end{tabular}}\end{table}

\begin{table}[t]\centering\scriptsize
\caption{Clinical secondary metrics on the unaltered test sets (Table~\ref{m:tab:natural}): mean (SD) over training seeds. Sensitivity and specificity at the operating point fixed on clean validation data; ECE with 10 equal-width bins; AUPRC baseline = prevalence.}\label{tab:natural_clinical}
\begin{tabular}{llcccccc}\toprule
Test set & Arm & AUROC & AUPRC & Brier & ECE & Sens & Spec \\\midrule
Thyroid (official split) ($\pi=0.38$) & ERM & 0.736 (0.009) & 0.606 (0.012) & 0.207 (0.005) & 0.082 (0.008) & 0.734 (0.047) & 0.619 (0.061) \\
 & mask & 0.731 (0.006) & 0.632 (0.007) & 0.214 (0.006) & 0.120 (0.026) & 0.347 (0.095) & 0.888 (0.041) \\
 & U-MtE & 0.742 (0.003) & 0.638 (0.007) & 0.203 (0.001) & 0.085 (0.007) & 0.425 (0.088) & 0.854 (0.050) \\
 & balanced & 0.716 (0.007) & 0.584 (0.012) & 0.219 (0.004) & 0.110 (0.007) & 0.714 (0.052) & 0.599 (0.061) \\
 & U-MtE$_{\text{bal}}$ & 0.731 (0.004) & 0.630 (0.007) & 0.207 (0.001) & 0.084 (0.006) & 0.399 (0.069) & 0.866 (0.034) \\
\midrule
ISIC, BCN held out ($\pi=0.23$) & ERM & 0.786 (0.003) & 0.586 (0.003) & 0.239 (0.006) & 0.279 (0.011) & 0.842 (0.038) & 0.519 (0.075) \\
 & mask & 0.769 (0.004) & 0.571 (0.009) & 0.286 (0.007) & 0.332 (0.009) & 0.860 (0.013) & 0.439 (0.023) \\
 & U-MtE & 0.773 (0.004) & 0.574 (0.008) & 0.280 (0.009) & 0.328 (0.012) & 0.858 (0.039) & 0.448 (0.089) \\
 & balanced & 0.781 (0.005) & 0.585 (0.009) & 0.228 (0.006) & 0.257 (0.011) & 0.811 (0.046) & 0.554 (0.087) \\
 & U-MtE$_{\text{bal}}$ & 0.774 (0.004) & 0.568 (0.008) & 0.279 (0.011) & 0.325 (0.016) & 0.860 (0.027) & 0.447 (0.051) \\
\midrule
ISIC, MSK held out ($\pi=0.19$) & ERM & 0.770 (0.008) & 0.477 (0.012) & 0.172 (0.007) & 0.161 (0.019) & 0.608 (0.049) & 0.786 (0.039) \\
 & mask & 0.790 (0.004) & 0.512 (0.009) & 0.146 (0.003) & 0.115 (0.009) & 0.593 (0.052) & 0.824 (0.037) \\
 & U-MtE & 0.783 (0.004) & 0.500 (0.008) & 0.150 (0.003) & 0.126 (0.009) & 0.566 (0.073) & 0.831 (0.049) \\
 & balanced & 0.766 (0.008) & 0.462 (0.014) & 0.178 (0.008) & 0.170 (0.020) & 0.626 (0.047) & 0.755 (0.046) \\
 & U-MtE$_{\text{bal}}$ & 0.774 (0.006) & 0.482 (0.009) & 0.154 (0.004) & 0.131 (0.010) & 0.594 (0.061) & 0.798 (0.051) \\
\midrule
ISIC, HAM held out ($\pi=0.11$) & ERM & 0.750 (0.005) & 0.353 (0.005) & 0.162 (0.006) & 0.234 (0.011) & 0.535 (0.092) & 0.813 (0.068) \\
 & mask & 0.789 (0.002) & 0.404 (0.006) & 0.143 (0.005) & 0.217 (0.011) & 0.559 (0.112) & 0.842 (0.063) \\
 & U-MtE & 0.799 (0.008) & 0.406 (0.009) & 0.145 (0.007) & 0.226 (0.016) & 0.573 (0.067) & 0.846 (0.041) \\
 & balanced & 0.747 (0.017) & 0.346 (0.019) & 0.156 (0.014) & 0.228 (0.015) & 0.539 (0.078) & 0.803 (0.076) \\
 & U-MtE$_{\text{bal}}$ & 0.802 (0.005) & 0.401 (0.009) & 0.139 (0.006) & 0.215 (0.021) & 0.606 (0.054) & 0.825 (0.037) \\
\midrule
Capsule (held-out frames) ($\pi=0.56$) & ERM & 0.889 (0.019) & 0.897 (0.046) & 0.130 (0.016) & 0.042 (0.011) & 0.770 (0.062) & 0.844 (0.055) \\
 & mask & 0.967 (0.011) & 0.971 (0.014) & 0.069 (0.013) & 0.029 (0.011) & 0.859 (0.032) & 0.933 (0.025) \\
 & U-MtE & 0.962 (0.009) & 0.967 (0.013) & 0.081 (0.017) & 0.065 (0.045) & 0.871 (0.042) & 0.897 (0.043) \\
 & balanced & 0.881 (0.024) & 0.889 (0.053) & 0.136 (0.017) & 0.051 (0.009) & 0.758 (0.087) & 0.838 (0.059) \\
 & U-MtE$_{\text{bal}}$ & 0.959 (0.008) & 0.964 (0.012) & 0.079 (0.007) & 0.044 (0.012) & 0.866 (0.048) & 0.897 (0.052) \\
\bottomrule
\end{tabular}\end{table}


\paragraph{Calibration} For the held-out BCN hospital masking worsens both Brier score and expected calibration error
(Table~\ref{tab:natural_clinical}); for HAM, MSK and capsule it improves every metric, in line with the AUROC results.
